\chapter{Method references and validation sources}
\label{app:references}

\begin{thebibliography}{99}
\bibitem{arr340}
Ak\c{c}elik, R., Besley, M. and Roper, R. (1999),
\emph{Fundamental Relationships for Traffic Flows at Signalised Intersections},
Research Report ARR 340, ARRB Transport Research Ltd, Vermont South, Australia.

\bibitem{akcelik1989}
Ak\c{c}elik, R. (1989),
\emph{On the Estimation of Lane Flows for Intersection Analysis},
Australian Road Research, 19(1), pp. 51--57; minor revisions 2001.

\bibitem{sidraShortLane}
SIDRA SOLUTIONS, \emph{Short Lane Model}, SIDRA INTERSECTION technical
documentation.  Consulted as a conceptual reference for short-lane interaction;
PySig does not claim to reproduce the proprietary SIDRA short-lane algorithm.

\bibitem{fhwaFlow}
Rouphail, N., Tarko, A. and Li, J.,
\emph{Traffic Flow at Signalized Intersections}, Traffic Flow Theory, Chapter 9,
Federal Highway Administration.

\bibitem{fhwaTiming}
Federal Highway Administration,
\emph{Traffic Signal Timing Manual}, Chapter 3: Operational and Safety Analysis.

\bibitem{rr67}
Transport and Road Research Laboratory, Research Report 67 (RR67), saturation-flow
relationships used as the basis for PySig's RR67-style unopposed estimator.
\end{thebibliography}

\section{Use of external methods}
Citation of a published method or commercial-model technical note does not imply that
PySig reproduces that complete method.  Where PySig implements an approximation, the
manual identifies the approximation explicitly.  This is particularly important for
short lanes, shared-lane blockage, priority movements, LOS and future ARR~340 queue-
discharge refinements.
